Sampling variation in a labour estimate contributes to its disagreement with a prediction. Compare frozen telemetry predictions by assigned intervention, role and task class on held-out teams or periods, while accounting for that sampling error. Correlation between binned averages leaves both systematic prediction error and individual prediction errors unresolved. The following identities show what a comparison against sampled labour can establish.

Treat each task's labour $H$ in a finite validation population as fixed. Let $f$ be its prediction, frozen before validation sampling, and suppose $\E_{\!s}[\widehat H]=H$ with finite design variance. Then
\begin{align*}
 \E_{\!s}[f-\widehat H]&=f-H,\\
 \E_{\!s}[(f-\widehat H)^2]&=(f-H)^2+\Var_{\!s}(\widehat H).
\end{align*}
Expand the square: the cross term vanishes because $\E_{\!s}[\widehat H-H]=0$. Dependence between tasks' samples does not alter these marginal identities, but affects aggregate uncertainty.

For a second prediction $g$ frozen before the same validation sample, subtraction cancels the common sampling variance:
\[
 \E_{\!s}[(f-\widehat H)^2-(g-\widehat H)^2]
 =(f-H)^2-(g-H)^2.
\]
Thus the observed squared-loss difference is design-unbiased for the true squared-loss difference. Estimating either loss level or the uncertainty of the difference requires further sampling-design information. Absolute error requires a separate analysis.

For example, eight one-hour person-intervals have allocations $(1,1,1,1,0,0,0,0)$. Sample one uniformly, with inclusion probability $1/8$. True labour is $H=4$, whereas its HT estimate is $0$ or $8$, each with probability $1/2$. A perfect prediction $f=4$ has apparent absolute error $4$ against the sampled estimate. The interval $[3.5,4.5]$ covers true labour surely but covers the sampled estimate with probability zero. The signed residual remains unbiased for the true prediction error. Absolute error and interval coverage against $\widehat H$ describe agreement with the noisy sampled total.

To assess accuracy against true labour, use a denser, independently adjudicated validation subset and report its remaining measurement error, or state a measurement model that identifies the chosen accuracy measure. Retain inclusion probabilities and joint sampling-design information for uncertainty calculations.

Validation should also test the score's sensitivity to task splitting, squashing and representation. Report reviewer disagreement and follow-up completeness alongside prediction error. To assess the consequence of those uncertainties, vary plausible missing labour and disputed retention scores. If those variations reverse an intervention ranking, the intervention decision depends on resolving or accepting the measurement uncertainty.
